% Generated draft: requires mathematical and source review.
\documentclass[11pt]{article}
\usepackage[margin=25mm]{geometry}
\usepackage{fontspec}
\setmainfont{Latin Modern Roman}
\usepackage{amsmath,amssymb,mathtools,unicode-math}
\setmathfont{Latin Modern Math}
\usepackage{enumitem,needspace,xurl,hyperref}
\setlength{\parindent}{0pt}
\setlength{\parskip}{5pt}
\setlength{\emergencystretch}{3em}
\newcommand{\N}{\mathbb{N}}
\newcommand{\Z}{\mathbb{Z}}
\newcommand{\Q}{\mathbb{Q}}
\newcommand{\R}{\mathbb{R}}
\newcommand{\C}{\mathbb{C}}
\newcommand{\F}{\mathbb{F}}
\newcommand{\A}{\mathbb{A}}
\newcommand{\PP}{\mathbb{P}}
\newcommand{\E}{\mathbb{E}}
\newcommand{\eps}{\varepsilon}
\newcommand{\dd}{\,\mathrm{d}}
\begin{document}
\section*{UnsolvedMath ID: 30001225}
\textbf{Source problem code:} catalogue entry
\subsection{Definitions and mathematical statement}
Let \(d,m,n\) be positive integers, and let \(\Omega\subset\C^d\) be a strictly convex domain with smooth boundary. Write
\[
L_a^2(\Omega)=\left\{f\in\mathcal O(\Omega):\int_\Omega |f(z)|^2\,\dd V(z)<\infty\right\},
\]
where \(\dd V\) is Lebesgue volume measure. The source uses the notation \(\mathcal A(\Omega)\), calling it the disk algebra. For a bounded domain, the usual interpretation is
\[
\mathcal A(\Omega)=\mathcal O(\Omega)\cap C(\overline\Omega),
\]
with the supremum norm. The source does not explicitly assume that \(\Omega\) is bounded or specify a convention for \(\mathcal A(\Omega)\) on an unbounded domain; this qualification remains unresolved.

Let \(A=(a_{ij})\) be an \(n\times m\) matrix with entries in the stated algebra \(\mathcal A(\Omega)\). Assume that there exists an integer \(r\), with \(0\le r\le\min\{m,n\}\), such that
\[
\operatorname{rank}A(\zeta)=r\qquad\text{for every }\zeta\in\partial\Omega.
\]
Thus the rank condition is imposed on the boundary, not throughout \(\Omega\). When the entries are bounded multipliers of the Bergman space, define
\[
M_A:L_a^2(\Omega)^m\longrightarrow L_a^2(\Omega)^n,
\qquad
(M_Af)_i=\sum_{j=1}^m a_{ij}f_j.
\]
In the usual bounded-domain interpretation above, this multiplication map is bounded. The source denotes its cokernel by
\[
F=\operatorname{coker}M_A.
\]
The algebraic quotient \(L_a^2(\Omega)^n/\operatorname{ran}M_A\) and the Hilbert-space quotient \(L_a^2(\Omega)^n/\overline{\operatorname{ran}M_A}\) agree as quotients only when the range is closed. The supplied record does not state the relevant closed-range conclusion or the complete privilege theorem, so the intended cokernel convention requires clarification.

\textbf{Requested classification.} For such domains and matrices, classify the analytic Hilbert modules \(F\) arising in the privilege theorem, and describe their parameter space within the framework of quasi-free Hilbert modules. The supplied original formulation presents these two targets as equivalent.

This is a classification request whose precise equivalence relation and parameter-space structure are not supplied. In particular, the record does not define the privilege condition, specify the intended definition of a quasi-free Hilbert module, explain whether the quotient itself must be quasi-free or is to be described through quasi-free presentations, or choose between unitary module equivalence and bounded invertible module isomorphism. These choices must be resolved before a unique precise classification problem can be stated. No additional interior rank condition, boundedness hypothesis, or equivalence relation is imposed here.
\subsection{Short English statement}
For a smooth strictly convex domain \(\Omega\subset\C^d\), the source asks for a classification and parameter space of privileged analytic Hilbert modules presented as cokernels of Bergman-space multiplication matrices whose entries lie in \(\mathcal A(\Omega)\) and whose boundary rank is constant. The privilege theorem, cokernel convention, equivalence relation, and role of quasi-free modules need clarification. The supplied open status is dated curation, not an independently verified conclusion.
\subsection{Sources}
\begin{itemize}
\item Supplied local UnsolvedMath registry record, permanent ID \(30001225\), source problem code \texttt{catalogue entry}, titled \emph{Classification of Analytic Hilbert Modules on Strictly Convex Domains}. The record belongs to the Oberwolfach Reports Open Problems collection and gives the proposed year as 2009. No individual proposer is supplied.
\item \emph{Hilbert Modules and Complex Geometry} (2009), Oberwolfach Reports, pp.~1047--1100, DOI: \texttt{10.4171/owr/2009/20}. Supplied source-report link and literature-source link labeled ``original report'': \url{https://doi.org/10.4171/owr/2009/20}.
\item The registry labels the statement \texttt{corrected\_verified} and reports that it was checked against the official Oberwolfach report. This is supplied source-verification metadata; the report has not been independently consulted for this draft.
\item Literature assessment dated \texttt{2026-08-21}: the supplied curation labels the problem open and describes the proposed classification beyond stated special cases as a broad open research problem. Those special cases are not specified in the supplied record. This assessment is supplied dated curation, not independently verified; the record explicitly calls for independent verification of time-sensitive status and cited claims.
\end{itemize}
\end{document}
